% GENERATED FILE. Edit canonical lessons, then run npm run generate:books.
% canonical-source-hash: fnv1a64:ce3564264903d33b
% canonical-lessons: AR-C473-waha, AR-C473-sarab, AR-C473-adghal, AR-C473-mustanqa, AR-C473-qahil

\chapter{Desert and jungle}
\label{ch:ar-desert-and-jungle}
\hlchaptermodality{473}

\begin{chapteropening}
\textbf{By the end of this chapter:} \emph{I can say an oasis, a mirage, a jungle, a swamp and barren.}

Complete the last lesson of chapter 473: I can say an oasis, a mirage, a jungle, a swamp and barren.
\end{chapteropening}

\section[\texorpdfstring{\ar{واحة} (wāḥa)}{واحة (wāḥa)}]{\ar{واحة} (wāḥa) — an oasis}

\label{lesson:AR-C473-waha}



\begin{quote}
Before the new one: say the Arabic for a slope, then the Arabic for an edge.
\end{quote}

\subsection*{You'll want to know: \ar{واحة}}
\textbf{\ar{واحة}} — \emph{wāḥa} — \textquotedblleft{}an oasis\textquotedblright{}.

\begin{quote}
\ar{نزلنا} \ar{في} \ar{واحة} \ar{صغيرة} \ar{في} \ar{الصحراء}.
\end{quote}

\begin{quote}
\emph{nazalnā fī wāḥa ṣaghīra fī aṣ-ṣaḥrāʾ.}
\end{quote}

\begin{quote}
We stopped at a small oasis in the desert.
\end{quote}

\begin{cousinweb}
\textbf{\ar{واحة}} (\emph{wāḥa}) was borrowed from Coptic, which goes back to ancient Egyptian (EtymDB); English \emph{oasis} comes from the same word.
\end{cousinweb}

\begin{grammarlens}[title={What you've built}]
One more word for getting around.
\end{grammarlens}

\subsection*{Your turn}
\begin{itemize}
\raggedright
  \item \emph{Say it:} \emph{wāḥa}
  \item \emph{Say it:} \emph{wāḥa}, once more
  \item \emph{Say it:} \emph{nazalnā fī wāḥa ṣaghīra fī aṣ-ṣaḥrāʾ.} — We stopped at a small oasis in the desert.
  \item \emph{Recall:} say the Arabic for a slope, then the Arabic for an edge, then say \emph{wāḥa} again
\end{itemize}

\begin{tcolorbox}[breakable,colback=teal!4,colframe=teal!35!black,title={Before you move on}]
What does \ar{واحة} mean? (\textquotedblleft{}An oasis\textquotedblright{}.) Say it once more.
\end{tcolorbox}
\section[\texorpdfstring{\ar{سراب} (sarāb)}{سراب (sarāb)}]{\ar{سراب} (sarāb) — a mirage}

\label{lesson:AR-C473-sarab}



\begin{quote}
Before the new one: say the Arabic for an edge, then the Arabic for an oasis.
\end{quote}

\subsection*{You'll want to know: \ar{سراب}}
\textbf{\ar{سراب}} — \emph{sarāb} — \textquotedblleft{}a mirage\textquotedblright{}.

\begin{quote}
\ar{رأى} \ar{المسافر} \ar{سرابا} \ar{في} \ar{الصحراء}.
\end{quote}

\begin{quote}
\emph{raʾā al-musāfir sarāban fī aṣ-ṣaḥrāʾ.}
\end{quote}

\begin{quote}
The traveller saw a mirage in the desert.
\end{quote}

\begin{cousinweb}
\textbf{\ar{سراب}} (\emph{sarāb}) is built on the native Arabic root \emph{s-r-b}, \textquotedblleft{}to flow; to go away\textquotedblright{} (Lane). A mirage shimmers like flowing water and then is gone.
\end{cousinweb}

\begin{grammarlens}[title={What you've built}]
One more word for getting around.
\end{grammarlens}

\subsection*{Your turn}
\begin{itemize}
\raggedright
  \item \emph{Say it:} \emph{sarāb}
  \item \emph{Say it:} \emph{sarāb}, once more
  \item \emph{Say it:} \emph{raʾā al-musāfir sarāban fī aṣ-ṣaḥrāʾ.} — The traveller saw a mirage in the desert.
  \item \emph{Recall:} say the Arabic for an edge, then the Arabic for an oasis, then say \emph{sarāb} again
\end{itemize}

\begin{tcolorbox}[breakable,colback=teal!4,colframe=teal!35!black,title={Before you move on}]
What does \ar{سراب} mean? (\textquotedblleft{}A mirage\textquotedblright{}.) Say it once more.
\end{tcolorbox}
\section[\texorpdfstring{\ar{أدغال} (adghāl)}{أدغال (adghāl)}]{\ar{أدغال} (adghāl) — a jungle; thick undergrowth}

\label{lesson:AR-C473-adghal}



\begin{quote}
Before the new one: say the Arabic for an oasis, then the Arabic for a mirage.
\end{quote}

\subsection*{You'll want to know: \ar{أدغال}}
\textbf{\ar{أدغال}} — \emph{adghāl} — \textquotedblleft{}a jungle; thick undergrowth\textquotedblright{}.

\begin{quote}
\ar{تعيش} \ar{النمور} \ar{في} \ar{الأدغال}.
\end{quote}

\begin{quote}
\emph{taʿīshu an-numūr fī al-adghāl.}
\end{quote}

\begin{quote}
Tigers live in the jungle.
\end{quote}

\begin{cousinweb}
\textbf{\ar{أدغال}} (\emph{adghāl}) is usually traced to the native Arabic root \emph{d-gh-l}, \textquotedblleft{}a thicket\textquotedblright{}. It is a plural in form.
\end{cousinweb}

\begin{grammarlens}[title={What you've built}]
One more word for getting around.
\end{grammarlens}

\subsection*{Your turn}
\begin{itemize}
\raggedright
  \item \emph{Say it:} \emph{adghāl}
  \item \emph{Say it:} \emph{adghāl}, once more
  \item \emph{Say it:} \emph{taʿīshu an-numūr fī al-adghāl.} — Tigers live in the jungle.
  \item \emph{Recall:} say the Arabic for an oasis, then the Arabic for a mirage, then say \emph{adghāl} again
\end{itemize}

\begin{tcolorbox}[breakable,colback=teal!4,colframe=teal!35!black,title={Before you move on}]
What does \ar{أدغال} mean? (\textquotedblleft{}A jungle; thick undergrowth\textquotedblright{}.) Say it once more.
\end{tcolorbox}
\section[\texorpdfstring{\ar{مستنقع} (mustanqaʿ)}{مستنقع (mustanqaʿ)}]{\ar{مستنقع} (mustanqaʿ) — a swamp; a marsh}

\label{lesson:AR-C473-mustanqa}



\begin{quote}
Before the new one: say the Arabic for a mirage, then the Arabic for a jungle.
\end{quote}

\subsection*{You'll want to know: \ar{مستنقع}}
\textbf{\ar{مستنقع}} — \emph{mustanqaʿ} — \textquotedblleft{}a swamp; a marsh\textquotedblright{}.

\begin{quote}
\ar{يعيش} \ar{الضفدع} \ar{في} \ar{المستنقع}.
\end{quote}

\begin{quote}
\emph{yaʿīshu aḍ-ḍifdaʿ fī al-mustanqaʿ.}
\end{quote}

\begin{quote}
The frog lives in the swamp.
\end{quote}

\begin{cousinweb}
\textbf{\ar{مستنقع}} (\emph{mustanqaʿ}) is built on the native Arabic root \emph{n-q-ʿ}, \textquotedblleft{}standing water\textquotedblright{} (Lane). It is the place-noun of Form X: where water stands.
\end{cousinweb}

\begin{grammarlens}[title={What you've built}]
One more word for getting around.
\end{grammarlens}

\subsection*{Your turn}
\begin{itemize}
\raggedright
  \item \emph{Say it:} \emph{mustanqaʿ}
  \item \emph{Say it:} \emph{mustanqaʿ}, once more
  \item \emph{Say it:} \emph{yaʿīshu aḍ-ḍifdaʿ fī al-mustanqaʿ.} — The frog lives in the swamp.
  \item \emph{Recall:} say the Arabic for a mirage, then the Arabic for a jungle, then say \emph{mustanqaʿ} again
\end{itemize}

\begin{tcolorbox}[breakable,colback=teal!4,colframe=teal!35!black,title={Before you move on}]
What does \ar{مستنقع} mean? (\textquotedblleft{}A swamp; a marsh\textquotedblright{}.) Say it once more.
\end{tcolorbox}
\section[\texorpdfstring{\ar{قاحل} (qāḥil)}{قاحل (qāḥil)}]{\ar{قاحل} (qāḥil) — barren}

\label{lesson:AR-C473-qahil}



\begin{quote}
Before the new one: say the Arabic for a jungle, then the Arabic for a swamp.
\end{quote}

\subsection*{You'll want to know: \ar{قاحل}}
\textbf{\ar{قاحل}} — \emph{qāḥil} — \textquotedblleft{}barren\textquotedblright{}.

\begin{quote}
\ar{كانت} \ar{الأرض} \ar{قاحلة} \ar{بلا} \ar{ماء}.
\end{quote}

\begin{quote}
\emph{kānat al-arḍ qāḥila bi-lā māʾ.}
\end{quote}

\begin{quote}
The land was barren, without water.
\end{quote}

\begin{cousinweb}
\textbf{\ar{قاحل}} (\emph{qāḥil}) is usually traced to the native Arabic root \emph{q-ḥ-l}, \textquotedblleft{}to dry up\textquotedblright{}.
\end{cousinweb}

\begin{grammarlens}[title={What you've built}]
One more word for getting around.
\end{grammarlens}

\subsection*{Your turn}
\begin{itemize}
\raggedright
  \item \emph{Say it:} \emph{qāḥil}
  \item \emph{Say it:} \emph{qāḥil}, once more
  \item \emph{Say it:} \emph{kānat al-arḍ qāḥila bi-lā māʾ.} — The land was barren, without water.
  \item \emph{Recall:} say the Arabic for a jungle, then the Arabic for a swamp, then say \emph{qāḥil} again
\end{itemize}

\begin{tcolorbox}[breakable,colback=teal!4,colframe=teal!35!black,title={Before you move on}]
What does \ar{قاحل} mean? (\textquotedblleft{}Barren\textquotedblright{}.) Say it once more.
\end{tcolorbox}
